\documentclass[11pt,a4paper]{article}

\usepackage[T1]{fontenc}
\usepackage[utf8]{inputenc}
\usepackage{setspace}

\setlength{\parskip}{0.8em}
\setlength{\parindent}{0pt}

% ----------------------------
% LANGUAGE SWITCH (EDIT THIS)
% ----------------------------
\newif\ifen
\entrue   % <-- set \entrue for English
%\enfalse % <-- set \enfalse for Polish

\begin{document}

% ============================
% ENGLISH VERSION
% ============================
\ifen

\begin{center}
\Large\textbf{Can isotopes help us better understand clouds?}
\end{center}

Clouds play a crucial role in the Earth's climate and water cycle.
They regulate the distribution of rainfall, influence weather patterns,
and affect how much solar energy is absorbed or reflected by the atmosphere.
Despite their importance, many processes occurring inside clouds remain
difficult to observe directly and are still poorly represented in weather
and climate models.

One possible way to study these processes is through stable water isotopes.
Water molecules are not all identical: some contain slightly heavier forms
of hydrogen or oxygen. These differences cause isotopes to behave differently
during evaporation, condensation, freezing, and melting. As a result,
the isotopic composition of atmospheric water carries information about the
processes that have shaped a cloud during its lifetime.

The aim of this project is to investigate whether isotopic measurements can
be used to better understand how clouds evolve. To achieve this, a new
computer model will be developed that simulates the growth and evolution of
individual cloud droplets and ice particles while simultaneously tracking
their isotopic composition. The model will be compared with measurements
collected during aircraft campaigns and laboratory experiments.

A key question addressed by the project is whether isotopic observations
contain enough information to explain the isotopic state of a cloud in a
physically consistent way. Answering this question will help determine which
cloud processes leave measurable isotopic signatures and whether isotopes can
provide information about cloud properties that are otherwise difficult to
observe.

The project combines modern cloud modelling with atmospheric isotope
research, two fields that have rarely been brought together at this level of
detail. The expected outcome is a new modelling framework for studying
isotopes in clouds and a better understanding of the information contained in
isotopic measurements. In the longer term, the results may contribute to
improving the representation of clouds in atmospheric models and advancing
our understanding of the Earth's water cycle.

\fi

% ============================
% POLISH VERSION
% ============================
\ifen\else

\begin{center}
\Large\textbf{Czy izotopy mogą pomóc lepiej zrozumieć chmury?}
\end{center}

Chmury odgrywają kluczową rolę w klimacie Ziemi i globalnym obiegu wody.
Wpływają na rozkład opadów, kształtują warunki pogodowe oraz decydują o tym,
ile energii słonecznej jest pochłaniane lub odbijane przez atmosferę.
Pomimo ich znaczenia wiele procesów zachodzących wewnątrz chmur pozostaje
trudnych do bezpośredniej obserwacji i nadal jest niedostatecznie opisanych
w modelach pogody i klimatu.

Jednym ze sposobów badania tych procesów jest wykorzystanie stabilnych
izotopów wody. Cząsteczki wody nie są identyczne – niektóre zawierają
nieco cięższe odmiany wodoru lub tlenu. Powoduje to, że podczas parowania,
skraplania, zamarzania i topnienia zachowują się one nieco inaczej.
W efekcie skład izotopowy wody atmosferycznej przechowuje informacje
o procesach, które zachodziły podczas powstawania i rozwoju chmury.

Celem projektu jest sprawdzenie, czy pomiary izotopowe mogą pomóc
w lepszym zrozumieniu ewolucji chmur. W tym celu zostanie opracowany nowy
model komputerowy opisujący wzrost i rozwój pojedynczych kropelek chmurowych
oraz kryształków lodu przy jednoczesnym śledzeniu ich składu izotopowego.
Wyniki symulacji będą porównywane z danymi pochodzącymi z kampanii lotniczych
oraz eksperymentów laboratoryjnych.

Kluczowym pytaniem projektu jest to, czy obserwacje izotopowe zawierają
wystarczającą ilość informacji, aby w spójny fizycznie sposób wyjaśnić
skład izotopowy rzeczywistych chmur. Odpowiedź na to pytanie pozwoli
określić, które procesy chmurowe pozostawiają mierzalne sygnały izotopowe
oraz czy izotopy mogą dostarczać informacji o właściwościach chmur,
których nie można łatwo obserwować innymi metodami.

Projekt łączy nowoczesne modelowanie chmur z badaniami izotopów atmosferycznych
– dwie dziedziny, które dotychczas rzadko były ze sobą łączone na tak
szczegółowym poziomie. Oczekiwanym rezultatem będzie nowe narzędzie do
badania izotopów w chmurach oraz lepsze zrozumienie informacji zawartych
w pomiarach izotopowych. W dłuższej perspektywie wyniki projektu mogą
przyczynić się do udoskonalenia modeli atmosferycznych oraz pogłębienia
wiedzy o obiegu wody w systemie klimatycznym Ziemi.

\fi

\end{document}